\section{Introduction}
\label{sec:intro}

This paper is a companion to the SVD cross-section compression paper.
The headline of that paper is a physics-and-correctness result; this
one is the engineering note that goes with it. After we shipped GPU
backends for SVD reconstruction, recursive-CSG geometry, photon
transport, and a multigroup random-ray solver, we found ourselves
re-learning a small set of lessons that the literature broadcasts but
that one only believes after living through them:

\begin{enumerate}
    \item Kernel-level optimisation matters less than data-flow
        amortisation. A per-XS-lookup speedup of $8$--$13\,\times$ on
        the SVD reconstruction microbench does not survive full
        transport on a 9-nuclide PWR pin cell; what does survive is
        whether the per-nuclide payloads (hundreds of MB, sometimes
        GB on actinide-heavy cases) are uploaded once or once per
        case.
    \item ``GPU is faster than CPU'' is a measurement-specific claim,
        not an architectural one. Our recursive-transport kernel
        runs $6.74\,\times$ faster than CPU \emph{on the constant-XS
        geometry-traversal microbench}; on full Godiva SVD transport
        with 20 rayon worker threads on a workstation Intel mobile,
        the same GPU kernel is $1.3\,\times$ \emph{slower}.
    \item Observability needs to be in the loop. We spent meaningful
        sessions trying to confirm by reading stdout that a cache
        ``actually fired''. The fix was a $\sim$200-line debug-metrics
        path that exposes cache entry counts, hit totals, byte
        usage, and the detected compute capability via PyO3, with a
        single CLI flag on the ICSBEP sweep driver that dumps one
        JSON line per case. Per-case deltas in the hit counter make
        cache behaviour falsifiable.
\end{enumerate}

The code is written in pure Rust against \texttt{cudarc} for the
CUDA driver and runtime APIs, and against NVRTC for at-startup
PTX compilation. There is no C dependency on the host side
either --- \texttt{hdf5-pure} reads the ENDF/B nuclear data
directly, and the same \texttt{Scene} JSON corpus drives both
backends.

\paragraph{Why CUDA.} Two reasons, in this order. (1) The amount
of available parallelism on a single device dwarfs what any
realistic host CPU offers. A consumer-tier RTX 3080 carries
$8\,704$ CUDA cores against the $20$ logical threads of the
workstation Intel mobile we ran on; a datacenter H100 carries
$16\,896$. Even after the launch-overhead and divergent-memory
penalties documented in §\ref{sec:honest} cut the effective ratio
by an order of magnitude, the device-side compute headroom is
where every scaled-up problem lives. (2) The vendor-specific
choice of CUDA over alternatives (HIP, SYCL, oneAPI) follows from
the deployment target: essentially all datacenter Monte Carlo
runs on NVIDIA hardware --- the clusters, the cloud spot pools,
the published OpenMC / Shift / MC21 GPU work. A portable layer
would be the right answer if the engine were targeting a
heterogeneous fleet. The engine targets the fleet it can
actually land on.

\paragraph{What this paper covers.} §\ref{sec:architecture}
describes the shared dispatch trait that lets CPU and GPU paths run
the same scene. §\ref{sec:caches} is the engineering core:
per-nuclide kernels, the new S($\alpha$,$\beta$) buffer cache, the
\texttt{Arc::as\_ptr} trick that makes both work, and what they
weigh in practice. §\ref{sec:sm_occupancy} catalogues the SM
occupancy and register-pressure findings, including the
\texttt{nuc\_t[128]} per-thread stack array that pins our PTX to
spillage on small cards. §\ref{sec:refill} reports the
PHYSOR~2022 Optimization~F refill-pool work with measured
saturation curves on RTX A1000 / RTX 3080.
§\ref{sec:adaptive_arch} describes the runtime
compute-capability detection that ended the
``rebuild-per-card'' tax. §\ref{sec:observability} introduces
\texttt{-{}-debug-metrics}. §\ref{sec:honest} consolidates the
``what we thought vs.\ reality'' findings; §\ref{sec:conclusion}
closes.

\paragraph{What this paper does \emph{not} cover.} The kernel-level
SVD reconstruction algorithm and its arithmetic-intensity tradeoffs
are documented in the companion SVD paper. Algorithm-level
correctness against ICSBEP handbook values, OpenMC cross-validation,
and the ENDF/B-VII.1 vs.\ VIII.1 library transition are documented
in the project's \texttt{STATUS.md} and benchmark CSVs; we cite a
few headline numbers but do not re-derive them here.
